\section{偏函数的限制和延拓}

\begin{definition}{偏函数在集合上一致}{}
	设$f$和$g$是偏函数,$E$是集合.若$E$包含于$f$和$g$各自的定义域,并且对任意$x\in E$有$f\left(x\right)=g\left(x\right)$,则称$f$和$g$在$E$上一致.
\end{definition}

\begin{proposition}{函数图像决定函数一致}{}
	如果偏函数$f$和$g$有相同的函数图像,那么$f$和$g$在$f$的定义域上一致.
\end{proposition}

\begin{definition}{函数的相等}{}
	设$f$和$g$是函数.若$f$和$g$的定义域都等于一个$A$,靶相同,则称$f$和$g$相等.
\end{definition}

\begin{proposition}{}{}
	设$f\coloneq\left(F,A,B\right)$和$g\coloneq\left(G,C,D\right)$是偏函数,则$F\subseteq G$ $\iff$ $A\subseteq C$且$f$和$g$在$A$上一致.
\end{proposition}

\begin{definition}{偏函数的延拓,子族}{}
	设$f\coloneq\left(F,A,B\right)$和$g\coloneq\left(G,C,D\right)$是偏函数.若$F\subseteq G$且$B\subseteq D$,则称$g$是$f$(到$C$)的延拓(或扩张).将$g$称为$D$的元素族时,称$f$是$g$的子族.
\end{definition}

\begin{definition}{偏函数的限制}{}
	设$f$是偏函数,陪域为$B$,$X$是$A$的定义域的子集,$F$为公式$x\in A\wedge y=f\left(x\right)$的图像.把$\left(F,X,B\right)$记作$f|_{X}$,称为$f$在$X$上的限制.
\end{definition}

\begin{proposition}{偏函数的限制的性质}{}
	\begin{enumerate}
		\item 每个函数都是它的每个限制的一个延拓.
		\item 如果偏函数$f,g$的靶相同,并且在集合$E$上一致,则$f|_{E}=g|_{E}$.
	\end{enumerate}
\end{proposition}













































































